\clearpage
\section{Scope and Goals}

The goal of this project is to design and implement a telemetry system capable of integrating with a variety of peripherals
used by a custom hydrogen-fueled vehicle.

The system is intended to collect operational data that allows the racing team to monitor lap count and lap times in real time,
while also providing insight into the system's behavior for troubleshooting and post-race analysis. By collecting this information
the telemetry system supports both performance evaluation and fault diagnosis.

The telemetry system is designed to be deployed on two different vehicles, each with its own specific
characteristics and requirements. As a result, the system must be modular in nature allowing individual
components to be adapted or replaced without requiring a complete redesign. This modular approach also
ensures that the system can be extended in the future should additional sensors, peripherals, or features be
required. Flexibility and scalability are therefore key design considerations throughout the project.

Another important objective is to simplify the development and debugging process. To achieve this, the
system must provide a convenient method for accessing logs and diagnostic information from the
target hardware. This capability is particularly important during testing and deployment, where limited
physical access and time constraints can make troubleshooting challenging. Reliable logging helps reduce
development time and improves overall system robustness.

Before development could begin, careful consideration was given to the selection of appropriate communication
technologies. Choosing the correct protocols is critical, as they directly affect system reliability, performance
and ease of integration with other vehicle electronics. Following the selection of the software and communication stack,
suitable hardware components were chosen to meet the project's functional and environmental requirements. The rationale
behind these decisions is discussed in detail in later chapters.

This thesis begins with an overview of the communication protocols selected for interfacing with the vehicle
and for transmitting the collected telemetry data. It then presents the hardware platform and explains the criteria
that guided its selection. Next, the overall system architecture is described, highlighting the interaction between
software and hardware components. Finally, the implementation details, testing results and conclusions are presented,
providing an evaluation of the system's performance and outlining potential areas for future improvement.